\section{Padronização das Placas}
\label{sec:padronizacao_placas}
Esta seção estabelece as diretrizes de projeto e as restrições metodológicas adotadas para o desenvolvimento das placas de aquisição e controle (SIBA e SAME). Visando otimizar os custos de manufatura e garantir a modularidade do sistema, por conta disso irá se usar placas de fenolite manufaturadas a mão, bem como módulos, placas de desenvolvimento e conversores comerciais.

Adotou-se o microcontrolador CH32V307VCT6 (WCH) como unidade de processamento
das placas SIBA e SAME. A escolha fundamenta-se em quatro critérios principais.
Primeiro, o conversor analógico-digital com taxa de até 1 Msps
atende à amostragem rápida de correntes e tensões exigida pela aquisição e pelo
controle. Segundo, a ampla memória disponível acomoda o \textit{firmware} sobre
o sistema operacional de tempo real Zephyr, as estruturas de dados de aquisição
e o processamento em ponto flutuante do controle vetorial e do observador, este
último amparado pela unidade de ponto flutuante (FPU) do núcleo RISC-V de
\SI{144}{\mega\hertz}. Terceiro, o dispositivo integra interface Ethernet
(MAC e PHY), detalhada adiante. Por fim, o baixo custo — cerca de
R\$30 por placa, ou aproximadamente R\$90 considerando
tributos de importação — viabiliza a produção de múltiplas unidades dentro das
restrições orçamentárias do projeto.

A interface Ethernet merece destaque por ter motivado a escolha. Durante o ano
de 2025, membros da equipe Albatroz Aerodesign identificaram a necessidade de um
método de comunicação mais simples e padronizado entre as placas de aquisição e
os computadores de ensaio. A conectividade por rede via cabo atende a essa
demanda: dispensa conversores de sinal intermediários, como adaptadores
UART-TTL (\textit{Universal Asynchronous Receiver/Transmitter} e \textit{Transistor-Transistor Logic}, e equipamentos dedicados de maior porte, uma vez que a maior parte dos
computadores e \textit{notebooks} dispõe nativamente de conexão de rede cabeada, além de facilidade adicionar futuras demandas não esperadas neste trabalho.

\subsection{Sistema operacional - Vcore}

O \textit{firmware} das duas placas desenvolvidas neste trabalho é estruturado
em três camadas lógicas, apresentadas na \autoref{fig:Metodologia_Diagrama_placas}.

\begin{figure}[H]
    \centering
    \caption{Camadas lógicas do \textit{firmware}}
    \label{fig:Metodologia_Diagrama_placas}
    \begin{tikzpicture}[
        font=\small,
        layer/.style={draw, rounded corners, minimum width=8cm, minimum height=1cm, align=center}
    ]
        \node[layer, fill=black!5] (app) {\textbf{Application}};
        \node[layer, below=0.8cm of app, fill=black!10] (sys) {\textbf{System}};
        \node[layer, below=0.8cm of sys, fill=black!15] (vcore) {\textbf{Vcore - OS}};

        \draw[<->, thick] (app.south) -- (sys.north);
        \draw[<->, thick] (sys.south) -- (vcore.north);
    \end{tikzpicture}

    Fonte: Autor (2026)
\end{figure}

A camada base (\textit{Vcore - OS}) encapsula as rotinas nativas do
\textit{Zephyr} e os periféricos do microcontrolador, apresentando uma
arquitetura idêntica para a SIBA e a SAME e promovendo o reaproveitamento de
código entre as duas placas. Ela é composta por cinco módulos, detalhados na
\autoref{fig:Metodologia_Vcore}:

\begin{itemize}
    \item \textbf{OSSI (\textit{Operating System Standard Interface}):} fila de trabalho única do sistema, que recebe os dados publicados pelo HAL e pelo CMD\_Core e despacha o processamento de lógica para a camada \textit{System} sempre na mesma \textit{thread} (\autoref{fig:arquitetura_fila_unica}).
    \item \textbf{NCI (\textit{Network Core Interface}):} gerenciador do \textit{socket} UDP, responsável pela recepção e pela transmissão de pacotes com latência minimizada.
    \item \textbf{HAL (\textit{Hardware Abstraction Layer}):} interface direta com os periféricos do microcontrolador (como o ADC operando via DMA para leituras analógicas em \textit{background}), isolando os registradores físicos das camadas superiores.
    \item \textbf{CMD\_Core:} interpretador central que recebe os pacotes binários brutos entregues pelo NCI, valida as instruções e as encaminha ao OSSI para despacho.
    \item \textbf{NVM (\textit{Non-Volatile Memory}):} guarda numa memória EEPROM externa, acessada por barramento I\textsuperscript{2}C, os coeficientes de calibração de cada sensor --- e, na SAME, também os parâmetros do observador. Na partida o módulo restaura o que está gravado; quando o operador confirma uma calibração pela aplicação de supervisão, o CMD\_Core manda gravá-la, de modo que ela sobrevive à desenergização da placa.
\end{itemize}

\begin{figure}[H]
    \centering
    \caption{Módulos da camada \textit{Vcore - OS}}
    \label{fig:Metodologia_Vcore}
    \begin{tikzpicture}[
        font=\small,
        box/.style={draw, rounded corners, align=center, minimum height=0.9cm, minimum width=2.4cm},
        arrow/.style={<->, thick}
    ]
        \node[box] (ossi) {OSSI};
        \node[box, above=1.1cm of ossi] (hal) {HAL};
        \node[box, left=1.6cm of ossi] (nci) {NCI};
        \node[box, right=1.6cm of ossi] (cmd) {CMD\_Core};
        \node[box, below=1.1cm of ossi] (nvm) {NVM\\(EEPROM)};

        \draw[arrow] (hal) -- (ossi);
        \draw[arrow] (nci) -- (ossi);
        \draw[arrow] (cmd) -- (ossi);
        \draw[arrow] (nvm) -- (ossi);
    \end{tikzpicture}

    Fonte: Autor (2026)
\end{figure}

Já a camada \textit{System} concentra os gerenciadores de domínio (bateria,
hélice, segurança, telemetria, ou o observador SMO no caso da SAME), sem
nenhuma \textit{thread} própria --- eles são chamados diretamente pela função
de despacho do OSSI, conforme detalhado adiante para cada placa
(\autoref{sec:siba} e \autoref{sec:same}). Por fim, a camada
\textit{Application} opera isolada no \textit{User Space}, restrita ao
barramento OSSI para qualquer interação com as camadas inferiores.

A alternativa mais direta ao problema de várias fontes de evento concorrentes
--- uma \textit{thread} própria por sensor, cada uma alterando diretamente as
variáveis compartilhadas --- foi descartada por expor o sistema a condições de
corrida e à necessidade de um \textit{mutex} por variável compartilhada, com o
consequente risco de \textit{deadlock} entre travas tomadas em ordens
diferentes por \textit{threads} diferentes. Adotou-se, em seu lugar, o padrão
de \textit{thread} única de processamento ilustrado na
\autoref{fig:arquitetura_fila_unica}: as únicas \textit{threads} que
permanecem bloqueando à espera de um evento externo são as que dependem
diretamente de periférico ou de rede (leitura de ADC, \textit{bit-banging} de
sensor, recepção UDP); ao terminarem, elas apenas empacotam o valor lido em
uma fila de dados (\texttt{k\_msgq}) e delegam todo o processamento de lógica
--- atualização de telemetria, verificação de segurança, despacho de comando
--- à fila de trabalho única (\texttt{k\_work\_q}), que executa sempre na
mesma \textit{thread}. Por construção, dois gerenciadores de domínio nunca
executam simultaneamente disputando o mesmo estado, o que dispensa \textit{mutex} entre
eles.

\begin{figure}[H]
    \centering
    \caption{Arquitetura de fila de trabalho única, comum às placas SIBA e SAME}
    \label{fig:arquitetura_fila_unica}
    \resizebox{\textwidth}{!}{%
    \begin{tikzpicture}[
        node distance=0.9cm and 1cm,
        font=\small,
        box/.style={draw, rounded corners, align=center, minimum height=0.9cm, minimum width=2.2cm},
        arrow/.style={-{Latex}, thick}
    ]
        \node[box] (hw) {Periférico\\(ADC, HX711, \dots)};
        \node[box, right=of hw] (thr) {\textit{Thread}\\bloqueante};
        \node[box, right=of thr] (msgq) {Fila de dados\\(\texttt{k\_msgq})};
        \node[box, right=of msgq] (wq) {Fila de trabalho\\única (\texttt{k\_work\_q})};
        \node[box, above right=0.5cm and 1cm of wq] (mgrs) {Gerenciadores\\de domínio};
        \node[box, below right=0.5cm and 1cm of wq] (nci) {NCI\\(rede UDP)};
        \node[box, right=of nci] (pc) {Aplicação\\de PC};
        \node[box, below=of msgq] (cmdq) {Fila de\\comandos};

        \draw[arrow] (hw) -- (thr);
        \draw[arrow] (thr) -- (msgq);
        \draw[arrow] (msgq) -- (wq);
        \draw[arrow] (wq) -- (mgrs);
        \draw[arrow] (wq) -- (nci);
        \draw[arrow] (nci) -- (pc);
        \draw[arrow] (pc) -- (nci);
        \draw[arrow] (nci) |- (cmdq);
        \draw[arrow] (cmdq) -- (wq);
    \end{tikzpicture}%
    }

    Fonte: Autor (2026)
\end{figure}
